\documentclass[10pt,a4paper]{amsart}
\usepackage{amsmath,amssymb,amsthm,graphicx}
% Retained minimal preamble for the Stacks Project chapter flat.tex (More on Flatness).
% Verbatim excerpt of the chapter preamble: Xy-pic setup (native preamble lines 14-18) and the native
% theorem declarations with the mathematical macros (native preamble lines 154-217). Copied mechanically;
% nothing added, removed or redefined. The wrapper class and the amsmath/amssymb/amsthm/graphicx base
% packages are supplied by the assembler, so the class/externaldocument part of the native preamble is omitted.
% For commutative diagrams we use Xy-pic
\usepackage[all]{xy}

% We use 2cell for 2-commutative diagrams.
\xyoption{2cell}
% Theorem environments.
%
\theoremstyle{plain}
\newtheorem{theorem}[subsection]{Theorem}
\newtheorem{proposition}[subsection]{Proposition}
\newtheorem{lemma}[subsection]{Lemma}

\theoremstyle{definition}
\newtheorem{definition}[subsection]{Definition}
\newtheorem{example}[subsection]{Example}
\newtheorem{exercise}[subsection]{Exercise}
\newtheorem{situation}[subsection]{Situation}

\theoremstyle{remark}
\newtheorem{remark}[subsection]{Remark}
\newtheorem{remarks}[subsection]{Remarks}

\numberwithin{equation}{subsection}

% Macros
%
\def\lim{\mathop{\mathrm{lim}}\nolimits}
\def\colim{\mathop{\mathrm{colim}}\nolimits}
\def\Spec{\mathop{\mathrm{Spec}}}
\def\Hom{\mathop{\mathrm{Hom}}\nolimits}
\def\Ext{\mathop{\mathrm{Ext}}\nolimits}
\def\SheafHom{\mathop{\mathcal{H}\!\mathit{om}}\nolimits}
\def\SheafExt{\mathop{\mathcal{E}\!\mathit{xt}}\nolimits}
\def\Sch{\mathit{Sch}}
\def\Mor{\mathop{\mathrm{Mor}}\nolimits}
\def\Ob{\mathop{\mathrm{Ob}}\nolimits}
\def\Sh{\mathop{\mathit{Sh}}\nolimits}
\def\NL{\mathop{N\!L}\nolimits}
\def\CH{\mathop{\mathrm{CH}}\nolimits}
\def\proetale{{pro\text{-}\acute{e}tale}}
\def\etale{{\acute{e}tale}}
\def\QCoh{\mathit{QCoh}}
\def\Ker{\mathop{\mathrm{Ker}}}
\def\Im{\mathop{\mathrm{Im}}}
\def\Coker{\mathop{\mathrm{Coker}}}
\def\Coim{\mathop{\mathrm{Coim}}}

% Boxtimes
%
\DeclareMathSymbol{\boxtimes}{\mathbin}{AMSa}{"02}

%
% Macros for moduli stacks/spaces
%
\def\QCohstack{\mathcal{QC}\!\mathit{oh}}
\def\Cohstack{\mathcal{C}\!\mathit{oh}}
\def\Spacesstack{\mathcal{S}\!\mathit{paces}}
\def\Quotfunctor{\mathrm{Quot}}
\def\Hilbfunctor{\mathrm{Hilb}}
\def\Curvesstack{\mathcal{C}\!\mathit{urves}}
\def\Polarizedstack{\mathcal{P}\!\mathit{olarized}}
\def\Complexesstack{\mathcal{C}\!\mathit{omplexes}}
% \Pic is the operator that assigns to X its picard group, usage \Pic(X)
% \Picardstack_{X/B} denotes the Picard stack of X over B
% \Picardfunctor_{X/B} denotes the Picard functor of X over B
\def\Pic{\mathop{\mathrm{Pic}}\nolimits}
\def\Picardstack{\mathcal{P}\!\mathit{ic}}
\def\Picardfunctor{\mathrm{Pic}}
\def\Deformationcategory{\mathcal{D}\!\mathit{ef}}

% --- appended 2026-10-12 (Chief): the project own definitions of the reference/slogan/history
% display environments, copied mechanically from the shared Stacks preamble lines 8-12
% (\usepackage{verbatim} plus the three \newenvironment lines, which the chapter preamble keeps
% verbatim from the project preamble). Nothing added, removed or redefined: these environments
% carry no mathematical content and the project itself typesets them as comments. Without these
% three lines a counted statement whose own text uses them fails with Environment slogan undefined.
\usepackage{verbatim}
\newenvironment{reference}{\comment}{\endcomment}
\newenvironment{slogan}{\comment}{\endcomment}
\newenvironment{history}{\comment}{\endcomment}
\usepackage{cleveref}
\crefname{originallemma}{original lemma}{original lemmas}
\Crefname{originallemma}{Original lemma}{Original lemmas}
\makeatletter
\newlabel{lemma-open-in-fibre-where-flat}{{05M9}{source}{}{sourcepointer-1}{}}
\newlabel{lemma-open-in-fibre-where-flat@cref}{{[originallemma][0][]05M9}{[1][0][]source}{}{}{}}
\makeatother
\crefname{originallemma}{original lemma}{original lemmas}
\Crefname{originallemma}{Original lemma}{Original lemmas}
\makeatletter
\newlabel{morphisms-lemma-base-change-module-flat}{{01U8}{source}{}{sourcepointer-2}{}}
\newlabel{morphisms-lemma-base-change-module-flat@cref}{{[originallemma][0][]01U8}{[1][0][]source}{}{}{}}
\makeatother
\crefname{originallemma}{original lemma}{original lemmas}
\Crefname{originallemma}{Original lemma}{Original lemmas}
\makeatletter
\newlabel{morphisms-lemma-dimension-fibre-after-base-change}{{02FY}{source}{}{sourcepointer-3}{}}
\newlabel{morphisms-lemma-dimension-fibre-after-base-change@cref}{{[originallemma][0][]02FY}{[1][0][]source}{}{}{}}
\makeatother
\crefname{originallemma}{original lemma}{original lemmas}
\Crefname{originallemma}{Original lemma}{Original lemmas}
\makeatletter
\newlabel{varieties-lemma-dimension-locally-algebraic}{{0A21}{source}{}{sourcepointer-4}{}}
\newlabel{varieties-lemma-dimension-locally-algebraic@cref}{{[originallemma][0][]0A21}{[1][0][]source}{}{}{}}
\makeatother
\title{Two equivalent formulations of the failure of flatness in dimensions below n}
\author{The Stacks Project Authors}
\date{stacks-project:flat.tex:local-snapshot-20261011; GNU FDL 1.2 or later}
\begin{document}
\maketitle
\noindent\textit{Editorial note.} All delivered mathematical text is the author's own, copied from the pinned Stacks Project chapter file flat.tex; no statement or proof step was rewritten, reordered or removed.\par\medskip
\noindent\textit{Editorial note.} Preamble provenance: the chapter's own retained minimal preamble (sources/stacks-flat/retained-preamble.tex, lines 1-85) is a verbatim excerpt of the native Xy-pic setup and of the native theorem declarations and mathematical macros, followed by the project's own reference/slogan/history display-environment definitions copied mechanically from the shared preamble; the wrapper class and the amsmath/amssymb/amsthm/graphicx base packages are supplied by the assembler, and extra\_packages is left empty because everything needed is already in the retained preamble. No macro was invented or redefined.\par\medskip
\noindent\textit{Editorial note.} Same-chapter pointers are located at their \textbackslash{}label lines inside the pinned chapter file; cross-chapter pointers are located in the project's own tag list staged in the snapshot (sources/stacks-chow/tags.txt), at the verbatim lines carrying their canonical tags.\par\medskip
\noindent\textit{Allowed standard prerequisites (editorial).}\begin{itemize}
\item Transcendence degree as the dimension of the closure of a point in the fibre, and the dimension formula for fibres after base change. Varieties, label varieties-lemma-dimension-locally-algebraic (proof lines 6290 and 6293); Morphisms, labels morphisms-lemma-base-change-module-flat and morphisms-lemma-dimension-fibre-after-base-change (proof lines 6296-6297). The first converts between the dimension of a closed subset of the fibre and the transcendence degree at its points, giving the equivalence even fibre by fibre, and the second two give the base change statement.
\item Openness of the flat locus inside a fibre. Same chapter flat.tex, label lemma-open-in-fibre-where-flat, invoked in the statement at proof line 6278. Makes the set of points of a fibre where the module is not flat a closed subset, so that its dimension is defined.
\end{itemize}\medskip
\ifdefined\hypertarget\hypertarget{sourcepointer-1}{}\fi
\noindent\textit{Source reference.} Original lemma 05M9: sources/stacks-flat/pinned-original-source.tex, source lines 2676--2676. Navigation pointer for \textbackslash{}ref\{lemma-open-in-fibre-where-flat\} used inside the counted span (line 6278); the argument that consumes it is retained in the span.\par\medskip
\ifdefined\hypertarget\hypertarget{sourcepointer-2}{}\fi
\noindent\textit{Source reference.} Original lemma 01U8: sources/stacks-chow/tags.txt, source lines 2401--2401. Navigation pointer for \textbackslash{}ref\{morphisms-lemma-base-change-module-flat\} used inside the counted span (line 6296); the argument that consumes it is retained in the span.\par\medskip
\ifdefined\hypertarget\hypertarget{sourcepointer-3}{}\fi
\noindent\textit{Source reference.} Original lemma 02FY: sources/stacks-chow/tags.txt, source lines 3209--3209. Navigation pointer for \textbackslash{}ref\{morphisms-lemma-dimension-fibre-after-base-change\} used inside the counted span (line 6297); the argument that consumes it is retained in the span.\par\medskip
\ifdefined\hypertarget\hypertarget{sourcepointer-4}{}\fi
\noindent\textit{Source reference.} Original lemma 0A21: sources/stacks-chow/tags.txt, source lines 12689--12689. Navigation pointer for \textbackslash{}ref\{varieties-lemma-dimension-locally-algebraic\} used inside the counted span (line 6290); the argument that consumes it is retained in the span.\par\medskip
\begin{lemma}
\label{lemma-pre-flat-dimension-n}
Let $f : X \to S$ be a morphism of schemes which is locally of finite type.
Let $\mathcal{F}$ be a quasi-coherent $\mathcal{O}_X$-module of finite
type. Let $n \geq 0$. The following are equivalent
\begin{enumerate}
\item for $s \in S$ the closed subset $Z \subset X_s$ of points
where $\mathcal{F}$ is not flat over $S$ (see
Lemma \ref{lemma-open-in-fibre-where-flat})
satisfies $\dim(Z) < n$, and
\item for $x \in X$ such that $\mathcal{F}$ is not flat at $x$
over $S$ we have $\text{trdeg}_{\kappa(f(x))}(\kappa(x)) < n$.
\end{enumerate}
If this is true, then it remains true after any base change.
\end{lemma}

\begin{proof}
Let $x \in X$ be a point over $s \in S$.
Then the dimension of the closure of $\{x\}$ in $X_s$
is $\text{trdeg}_{\kappa(s)}(\kappa(x))$ by
Varieties, Lemma \ref{varieties-lemma-dimension-locally-algebraic}.
Conversely, if $Z \subset X_s$ is a closed subset
of dimension $d$, then there exists a point $x \in Z$
with $\text{trdeg}_{\kappa(s)}(\kappa(x)) = d$ (same reference).
Therefore the equivalence of (1) and (2) holds (even fibre by fibre).
The statement on base change follows from
Morphisms, Lemmas \ref{morphisms-lemma-base-change-module-flat} and
\ref{morphisms-lemma-dimension-fibre-after-base-change}.
\end{proof}

\end{document}
